% ============================================================
\section{Operator Interface: Additive Horizon Systems}
\label{sec:bedrock}
% ============================================================

This section isolates the interface on which the structural calculus depends.
Only the algebraic and energetic data needed for projection identities, defect
identities, tower calculus, elimination, and solver constructions are fixed
here. No preferred coordinate model and no counting realization are imposed at
this stage.

% ------------------------------------------------------------
\subsection{Additive state spaces and energy}
\label{subsec:algebraic-foundations}
% ------------------------------------------------------------

\begin{definition}[Additive state space]\label{def:vector-space}\lean{AdditiveStateSpace}
An \emph{additive state space} is a set $\cH$ equipped with operations
\[
0 \in \cH,
\qquad
+ : \cH \times \cH \to \cH,
\qquad
- : \cH \to \cH,
\]
such that for all $x,y,z \in \cH$,
\begin{enumerate}[label=(\roman*), itemsep=0.2em]
\item $(x+y)+z = x+(y+z)$,
\item $x+y = y+x$,
\item $x+0 = x$,
\item $x+(-x)=0$.
\end{enumerate}
\end{definition}

\begin{definition}[Energy pairing]\label{def:inner-product}\lean{EnergyPairing}
An \emph{energy pairing} on an additive state space $\cH$ is a symmetric
positive bilinear form
\[
\langle \cdot,\cdot\rangle : \cH \times \cH \to \bbR
\]
such that
\[
\langle x,x\rangle \ge 0
\qquad\text{for all }x\in\cH,
\]
and $\langle x,x\rangle = 0$ only when $x=0$.
The induced energy norm is
\[
\|x\|^2 := \langle x,x\rangle.
\]
\end{definition}

\begin{definition}[Finite coherence space]\label{def:hilbert-space}\lean{FiniteCoherenceSpace}
A \emph{finite coherence space} is an additive state space equipped with an
energy pairing. Whenever orthogonality, adjoints, or defect energies are used
in later chapters, they are understood relative to that fixed pairing.
\end{definition}

% ------------------------------------------------------------
\subsection{Horizon towers and observed operators}

\begin{definition}[Horizon tower]
\label{def:horizon-tower}\lean{HorizonTower}
A \emph{horizon tower} on $\cH$ is a nested family of idempotent observable
projections $(\CCPi{h})_{h\in\bbN}$.
Nestedness means that coarse observation after finer observation changes
nothing:
\[
\CCPi{h}\CCPi{h'}=\CCPi{h}
\qquad (h\le h').
\]
\end{definition}

\begin{definition}[Shadow projector]\label{def:shadow-projector}\lean{shadowProj}
The shadow at horizon $h$ is the complementary operator
\[
\CCPibar{h}:=\id-\CCPi{h}.
\]
Thus every state splits into visible part plus shadow:
\[
x=\CCPi{h}x+\CCPibar{h}x.
\]
\end{definition}

\begin{definition}[Transitivity of horizons]\label{def:transitivity}\lean{nested_proj_comp}\lean{nested_proj_comp_ge}
For a nested horizon tower,
\[
\CCPi{h}\CCPi{h'}=\CCPi{h}
\qquad (h\le h'),
\qquad
\CCPi{h'}\CCPi{h}=\CCPi{h}
\qquad (h\le h').
\]
This is the algebraic coherence law for observation across different
resolutions.
\end{definition}

\begin{definition}[Energy-orthogonal tower]\label{def:orthogonal-tower}\lean{EnergyOrthogonalTower}
Let $\cH$ carry an energy pairing. A horizon tower is
\emph{energy-orthogonal} if each projection $\CCPi{h}$ is self-adjoint for that
pairing and
\[
\langle \CCPi{h}x,\CCPibar{h}y\rangle = 0
\qquad
\text{for all }x,y\in\cH.
\]
Equivalently, each horizon cut splits the state space into an orthogonal
observable sector and an orthogonal shadow sector.
\end{definition}

\begin{definition}[Additive operator]\label{def:linear-operator}\lean{AddMap}
An \emph{additive operator} on $\cH$ is a map $A:\cH\to\cH$ preserving addition
and zero:
\[
A(x+y)=Ax+Ay,
\qquad
A(0)=0.
\]
\end{definition}

\begin{definition}[Energy adjoint]\label{def:adjoint}\lean{HasEnergyAdjoint}\lean{defectEnergy_operator_form}
Given additive operators $A$ and $A^\ast$, one says that $A^\ast$ is an
\emph{energy adjoint} of $A$ if
\[
\langle Ax,y\rangle = \langle x,A^\ast y\rangle
\qquad
\text{for all }x,y\in\cH.
\]
This is exactly the adjoint structure needed for the operator form of defect
energy.
\end{definition}

\begin{definition}[Graded differential and nilpotency]
\label{def:nilpotent-differential}\lean{isNilpotent}
A \emph{nilpotent differential} on $\cH$ is an additive operator $d$ satisfying
\[
d^2=0.
\]
\end{definition}

\begin{definition}[Observed operator]
\label{def:observed-operator}\lean{observedOp}
Given a horizon tower and an additive operator $A$, the observed operator at
horizon $h$ is
\[
\CCObs{A}{h}:=\CCPi{h}A\CCPi{h}.
\]
\end{definition}

\begin{definition}[Leakage operator]
\label{def:leakage-operator}\lean{leakageOp}
The leakage operator is
\[
\CCLeak{A}{h}:=\CCPibar{h}A\CCPi{h}.
\]
It measures what leaves the observable sector when $A$ acts on visible input.
\end{definition}

\begin{definition}[Return operator]
\label{def:return-operator}\lean{returnOp}
The return operator is
\[
R_h:=\CCPi{h}A\CCPibar{h}.
\]
It measures what returns from shadow to the observable sector.
\end{definition}

\begin{remark}[Defect as round trip]
\label{rem:defect-round-trip}\lean{defect_factorization}
The defect is the observable round trip through shadow:
\[
\CCDefect{A}{h}=R_h\,\CCLeak{A}{h}
=\CCPi{h}A\CCPibar{h}A\CCPi{h}.
\]
Leakage and defect are therefore different objects. Leakage is escape;
defect is returned nonclosure.
\end{remark}

\paragraph{Interface principle.}
The structural theorems of Part~II depend only on the data above together with
the additional hypotheses stated locally in each chapter. No coordinate
realization, counting model, or physical interpretation is used at this stage.

\paragraph{Realizations enter later.}
Part~I proves that Closure Balance supplies a distinguished finite
realization of this interface. Chapter~13 proves that conservative
characteristic transport can furnish another realization. Both belong to the
realized-model layer, not to the structural interface itself.
